\documentclass{article}
\usepackage{amsmath, amssymb, amsthm}
\usepackage{geometry}
\usepackage{algorithm}
\usepackage{algpseudocode}
\geometry{margin=1in}

\title{Convergence Analysis of Clipped Adam}
\author{}
\date{}

\begin{document}
\maketitle

\section*{Algorithm Details Expansion}

\subsection*{Algorithm Name}
Clipped Adam (Adam with Gradient Clipping)

\subsection*{Mathematical Formulation}
Let $f: \mathbb{R}^d \to \mathbb{R}$ be the objective function. At each iteration $t$, the algorithm computes the gradient $g_t = \nabla f(x_t)$. To prevent excessively large updates, the gradient is scaled down if its $\ell_2$ norm exceeds a threshold $C > 0$. The clipped gradient is denoted as $\hat{g}_t$:
\begin{equation}
    \hat{g}_t = \frac{g_t}{\max\left(1, \frac{\|g_t\|}{C}\right)}
\end{equation}
The algorithm then maintains an exponential moving average of the clipped gradients (first moment, $m_t$) and their element-wise squares (second moment, $v_t$):
\begin{align}
    m_t &= \beta_1 m_{t-1} + (1 - \beta_1) \hat{g}_t \\
    v_t &= \beta_2 v_{t-1} + (1 - \beta_2) \hat{g}_t^2
\end{align}
where $\beta_1, \beta_2 \in [0, 1)$ are the decay rates. For theoretical clarity in the convergence proof, we analyze the standard uncorrected moments. The parameter vector $x_t$ is updated using a step size $\eta$ and a small constant $\epsilon > 0$ for numerical stability:
\begin{equation}
    x_{t+1} = x_t - \eta \frac{m_t}{\sqrt{v_t} + \epsilon}
\end{equation}
where operations in the denominator are element-wise. 

\subsection*{Pseudocode}
\begin{algorithm}[H]
\caption{Clipped Adam}
\begin{algorithmic}[1]
\State \textbf{Input:} Initial parameters $x_0 \in \mathbb{R}^d$, step size $\eta$, clipping threshold $C$, decay rates $\beta_1, \beta_2$, constant $\epsilon$
\State \textbf{Initialize:} $m_0 \gets 0$, $v_0 \gets 0$
\For{$t = 1$ to $T$}
    \State $g_t \gets \nabla f(x_{t-1})$
    \State $\hat{g}_t \gets g_t / \max(1, \|g_t\|/C)$
    \State $m_t \gets \beta_1 m_{t-1} + (1 - \beta_1) \hat{g}_t$
    \State $v_t \gets \beta_2 v_{t-1} + (1 - \beta_2) \hat{g}_t^2$
    \State $x_t \gets x_{t-1} - \eta \frac{m_t}{\sqrt{v_t} + \epsilon}$
\EndFor
\State \textbf{Return} $x_T$
\end{algorithmic}
\end{algorithm}

\subsection*{Dry Run Example}
Let $f(w) = (w-3)^2$. We initialize $w_0 = 0$. 
Hyperparameters: $\eta = 0.1$, $C = 1$, $\beta_1 = 0.9$, $\beta_2 = 0.99$, $\epsilon = 0.1$.

\textbf{Iteration 1 ($t=1$):}
\begin{itemize}
    \item Gradient: $g_1 = 2(w_0 - 3) = 2(0 - 3) = -6$
    \item Clip: $\|g_1\| = 6 > 1$, so $\hat{g}_1 = -6 / 6 = -1$
    \item Moments: $m_1 = 0.9(0) + 0.1(-1) = -0.1$
    \item Variance: $v_1 = 0.99(0) + 0.01(-1)^2 = 0.01$
    \item Update: $w_1 = 0 - 0.1 \frac{-0.1}{\sqrt{0.01} + 0.1} = 0 - 0.1 \frac{-0.1}{0.1 + 0.1} = 0.05$
\end{itemize}

\textbf{Iteration 2 ($t=2$):}
\begin{itemize}
    \item Gradient: $g_2 = 2(0.05 - 3) = -5.9$
    \item Clip: $\|g_2\| = 5.9 > 1$, so $\hat{g}_2 = -5.9 / 5.9 = -1$
    \item Moments: $m_2 = 0.9(-0.1) + 0.1(-1) = -0.19$
    \item Variance: $v_2 = 0.99(0.01) + 0.01(-1)^2 = 0.0099 + 0.01 = 0.0199$
    \item Update: $w_2 = 0.05 - 0.1 \frac{-0.19}{\sqrt{0.0199} + 0.1} \approx 0.05 + \frac{0.019}{0.1411 + 0.1} \approx 0.05 + 0.0788 = 0.1288$
\end{itemize}

\textbf{Iteration 3 ($t=3$):}
\begin{itemize}
    \item Gradient: $g_3 = 2(0.1288 - 3) = -5.7424$
    \item Clip: $\|g_3\| = 5.7424 > 1$, so $\hat{g}_3 = -1$
    \item Moments: $m_3 = 0.9(-0.19) + 0.1(-1) = -0.271$
    \item Variance: $v_3 = 0.99(0.0199) + 0.01(-1)^2 = 0.019701 + 0.01 = 0.029701$
    \item Update: $w_3 = 0.1288 - 0.1 \frac{-0.271}{\sqrt{0.029701} + 0.1} \approx 0.1288 + \frac{0.0271}{0.1723 + 0.1} \approx 0.1288 + 0.0995 = 0.2283$
\end{itemize}

\section*{Convergence Theory Development}

\subsection*{Assumptions}
We make the following standard assumptions for the non-convex optimization setting:
\begin{enumerate}
    \item \textbf{Smoothness:} The objective function $f$ is continuously differentiable and $L$-smooth, meaning there exists a constant $L > 0$ such that for all $x, y \in \mathbb{R}^d$:
    \begin{equation}
        f(y) \leq f(x) + \langle \nabla f(x), y - x \rangle + \frac{L}{2} \|y - x\|^2
    \end{equation}
    \item \textbf{Bounded Below:} The objective function is bounded below, meaning $f(x) \geq f^*$ for all $x \in \mathbb{R}^d$.
    \item \textbf{Bounded Gradients:} The gradients are globally bounded by some constant $G > 0$, such that $\|\nabla f(x)\| \leq G$ for all $x \in \mathbb{R}^d$.
\end{enumerate}

\subsection*{Main Convergence Theorem}
Under the stated assumptions, let $x_t$ be the sequence generated by Clipped Adam. Let $\Phi_t = \min(\|\nabla f(x_t)\|^2, C \|\nabla f(x_t)\|)$. Then, the average gradient measure satisfies:
\begin{equation}
    \frac{1}{T} \sum_{t=1}^T \Phi_t \leq \frac{f(x_1) - f^*}{T \eta \eta_L} + \frac{K}{\eta \eta_L}
\end{equation}
where $\eta_L = \frac{1}{C + \epsilon}$ and $K = O(\eta^2) + O(\eta \sqrt{1-\beta_2})$ is a constant encapsulating the errors strictly from momentum and adaptive learning rate variations. By choosing $\eta = O(1/\sqrt{T})$ and $1-\beta_2 = O(1/T)$, the algorithm converges to a stationary point at a rate of $O(1/\sqrt{T})$.

\subsection*{Theoretical Properties}
\begin{itemize}
    \item \textbf{Stability via Clipping:} In standard Adam, the effective step size can explode if $v_t$ is small while $m_t$ is large. Gradient clipping ensures that $\|\hat{g}_t\| \leq C$, which deterministically bounds the preconditioner and the momentum term, guaranteeing strictly bounded step sizes.
    \item \textbf{Momentum Handling:} The proof employs a virtual sequence $z_t$ which elegantly decouples the momentum history from the current gradient step, a critical property for proving convergence of momentum-based methods in non-convex settings.
\end{itemize}

\subsection*{Discussion}
The theorem shows that Clipped Adam converges to a region where $\min(\|\nabla f(x)\|^2, C \|\nabla f(x)\|)$ is small. For gradients larger than the clipping threshold $C$, the algorithm minimizes the gradient norm (acting like normalized SGD). For gradients smaller than $C$, it behaves like standard Adam and minimizes the squared gradient norm, successfully locating stationary points for non-convex functions.

\section*{Complete Convergence Proof}

\subsection*{Preliminaries (Definitions and Lemmas)}

Let us define the adaptive learning rate matrix at time $t$ as $H_t = \text{diag}\left(\frac{1}{\sqrt{v_t} + \epsilon}\right)$. The update rule is therefore $x_{t+1} = x_t - \eta H_t m_t$.

\textbf{Lemma 1 (Bounded Moments):} Due to gradient clipping, $\|\hat{g}_t\| \leq C$. By induction, since $m_t$ is a convex combination of $\hat{g}_t$ and $m_{t-1}$ (and $m_0=0$), we have $\|m_t\| \leq C$. Similarly, $v_{t,i} \leq C^2$.

\textbf{Lemma 2 (Bounded Preconditioner):} Since $0 \leq v_{t,i} \leq C^2$, the eigenvalues of $H_t$ are bounded. Let $\eta_L = \frac{1}{C+\epsilon}$ and $\eta_U = \frac{1}{\epsilon}$. We have $\eta_L I \preceq H_t \preceq \eta_U I$.

\textbf{Lemma 3 (Preconditioner Stability):} The change in the preconditioner is bounded.
\begin{proof}
    For any coordinate $i$, $v_{t,i} - v_{t-1,i} = (1-\beta_2)(\hat{g}_{t,i}^2 - v_{t-1,i})$. Since $\hat{g}_{t,i}^2 \leq C^2$ and $v_{t-1,i} \leq C^2$, we have $|v_{t,i} - v_{t-1,i}| \leq (1-\beta_2)C^2$. 
    Using the inequality $|\frac{1}{a+\epsilon} - \frac{1}{b+\epsilon}| \leq \frac{|a-b|}{\epsilon^2}$ for $a,b \geq 0$, and setting $a=\sqrt{v_{t,i}}, b=\sqrt{v_{t-1,i}}$:
    \begin{equation}
        |H_{t,i} - H_{t-1,i}| \leq \frac{|\sqrt{v_{t,i}} - \sqrt{v_{t-1,i}}|}{\epsilon^2} \leq \frac{\sqrt{|v_{t,i} - v_{t-1,i}|}}{\epsilon^2} \leq \frac{\sqrt{1-\beta_2} C}{\epsilon^2}
    \end{equation}
    Thus, the operator norm $\|H_t - H_{t-1}\|_2 \leq \frac{\sqrt{1-\beta_2} C}{\epsilon^2}$.
\end{proof}

\subsection*{Main Proof (Step-by-Step Derivation)}

\textbf{[Step 1]} We define a virtual sequence $z_t$ to handle the momentum term:
\begin{equation}
    z_t = x_t - \frac{\eta \beta_1}{1-\beta_1} H_{t-1} m_{t-1}
\end{equation}
Note that the distance between the virtual sequence and the actual sequence is bounded:
\begin{equation}
    \|x_t - z_t\| = \left\| \frac{\eta \beta_1}{1-\beta_1} H_{t-1} m_{t-1} \right\| \leq \frac{\eta \beta_1}{1-\beta_1} \|H_{t-1}\|_2 \|m_{t-1}\| \leq \frac{\eta \beta_1 \eta_U C}{1-\beta_1} \label{eq:dist_xz}
\end{equation}

\textbf{[Step 2]} We compute the difference $z_{t+1} - z_t$:
\begin{align}
    z_{t+1} - z_t &= x_{t+1} - x_t - \frac{\eta \beta_1}{1-\beta_1} H_t m_t + \frac{\eta \beta_1}{1-\beta_1} H_{t-1} m_{t-1} \nonumber \\
    &= - \eta H_t m_t - \frac{\eta \beta_1}{1-\beta_1} H_t m_t + \frac{\eta \beta_1}{1-\beta_1} H_{t-1} m_{t-1} \nonumber \\
    &= - \frac{\eta}{1-\beta_1} H_t m_t + \frac{\eta \beta_1}{1-\beta_1} H_{t-1} m_{t-1}
\end{align}

\textbf{[Step 3]} We substitute $m_t = \beta_1 m_{t-1} + (1-\beta_1) \hat{g}_t$ into the equation from [Step 2]:
\begin{align}
    z_{t+1} - z_t &= - \frac{\eta}{1-\beta_1} H_t (\beta_1 m_{t-1} + (1-\beta_1) \hat{g}_t) + \frac{\eta \beta_1}{1-\beta_1} H_{t-1} m_{t-1} \nonumber \\
    &= - \eta H_t \hat{g}_t + \frac{\eta \beta_1}{1-\beta_1} (H_{t-1} - H_t) m_{t-1}
\end{align}

\textbf{[Step 4]} We bound the norm of the step $\|z_{t+1} - z_t\|$ using Lemma 1, Lemma 2, and Lemma 3:
\begin{align}
    \|z_{t+1} - z_t\| &\leq \eta \|H_t\|_2 \|\hat{g}_t\| + \frac{\eta \beta_1}{1-\beta_1} \|H_{t-1} - H_t\|_2 \|m_{t-1}\| \nonumber \\
    &\leq \eta \eta_U C + \frac{\eta \beta_1}{1-\beta_1} \frac{\sqrt{1-\beta_2} C^2}{\epsilon^2} := \Delta_z
\end{align}

\textbf{[Step 5]} We apply the $L$-smoothness assumption on $f$ at the virtual point $z_t$:
\begin{equation}
    f(z_{t+1}) \leq f(z_t) + \langle \nabla f(z_t), z_{t+1} - z_t \rangle + \frac{L}{2} \|z_{t+1} - z_t\|^2
\end{equation}

\textbf{[Step 6]} We expand the inner product $\langle \nabla f(z_t), z_{t+1} - z_t \rangle$ using the identity from [Step 3]:
\begin{equation}
    \langle \nabla f(z_t), z_{t+1} - z_t \rangle = - \eta \langle \nabla f(z_t), H_t \hat{g}_t \rangle + \frac{\eta \beta_1}{1-\beta_1} \langle \nabla f(z_t), (H_{t-1} - H_t) m_{t-1} \rangle
\end{equation}
We bound the second term using Cauchy-Schwarz and the global gradient bound $G$:
\begin{equation}
    \frac{\eta \beta_1}{1-\beta_1} \langle \nabla f(z_t), (H_{t-1} - H_t) m_{t-1} \rangle \leq \frac{\eta \beta_1 G}{1-\beta_1} \frac{\sqrt{1-\beta_2} C^2}{\epsilon^2}
\end{equation}

\textbf{[Step 7]} We relate the first term $- \eta \langle \nabla f(z_t), H_t \hat{g}_t \rangle$ back to the exact gradient $\nabla f(x_t)$:
\begin{equation}
    - \eta \langle \nabla f(z_t), H_t \hat{g}_t \rangle = - \eta \langle \nabla f(x_t), H_t \hat{g}_t \rangle + \eta \langle \nabla f(x_t) - \nabla f(z_t), H_t \hat{g}_t \rangle
\end{equation}
By Cauchy-Schwarz and $L$-smoothness ($\|\nabla f(x_t) - \nabla f(z_t)\| \leq L \|x_t - z_t\|$), and using \eqref{eq:dist_xz}:
\begin{align}
    \eta \langle \nabla f(x_t) - \nabla f(z_t), H_t \hat{g}_t \rangle &\leq \eta L \|x_t - z_t\| \|H_t \hat{g}_t\| \nonumber \\
    &\leq \eta L \left( \frac{\eta \beta_1 \eta_U C}{1-\beta_1} \right) (\eta_U C) = \frac{L \eta^2 \beta_1 \eta_U^2 C^2}{1-\beta_1}
\end{align}

\textbf{[Step 8]} We analyze the main descent term $- \eta \langle \nabla f(x_t), H_t \hat{g}_t \rangle$. Since $H_t$ is a positive diagonal matrix with minimal element $\eta_L$, and substituting $\hat{g}_t = \nabla f(x_t) / \max(1, \|\nabla f(x_t)\|/C)$:
\begin{align}
    - \eta \langle \nabla f(x_t), H_t \hat{g}_t \rangle &\leq - \eta \eta_L \langle \nabla f(x_t), \hat{g}_t \rangle \nonumber \\
    &= - \eta \eta_L \frac{\|\nabla f(x_t)\|^2}{\max(1, \|\nabla f(x_t)\|/C)}
\end{align}
Notice that for any vector $v$, $\frac{\|v\|^2}{\max(1, \|v\|/C)} \geq \min(\|v\|^2, C\|v\|)$. Let $\Phi_t = \min(\|\nabla f(x_t)\|^2, C\|\nabla f(x_t)\|)$. Thus:
\begin{equation}
    - \eta \langle \nabla f(x_t), H_t \hat{g}_t \rangle \leq - \eta \eta_L \Phi_t
\end{equation}

\textbf{[Step 9]} We combine all bounds into the descent inequality from [Step 5]:
\begin{equation}
    f(z_{t+1}) \leq f(z_t) - \eta \eta_L \Phi_t + \frac{L \eta^2 \beta_1 \eta_U^2 C^2}{1-\beta_1} + \frac{\eta \beta_1 G \sqrt{1-\beta_2} C^2}{(1-\beta_1)\epsilon^2} + \frac{L}{2} \Delta_z^2
\end{equation}
Let $K$ be the sum of all constant error terms:
\begin{equation}
    K = \frac{L \eta^2 \beta_1 \eta_U^2 C^2}{1-\beta_1} + \frac{\eta \beta_1 G \sqrt{1-\beta_2} C^2}{(1-\beta_1)\epsilon^2} + \frac{L}{2} \left( \eta \eta_U C + \frac{\eta \beta_1 \sqrt{1-\beta_2} C^2}{(1-\beta_1)\epsilon^2} \right)^2
\end{equation}
Thus, $f(z_{t+1}) \leq f(z_t) - \eta \eta_L \Phi_t + K$.

\textbf{[Step 10]} We telescope the sum from $t=1$ to $T$:
\begin{equation}
    \eta \eta_L \sum_{t=1}^T \Phi_t \leq f(z_1) - f(z_{T+1}) + T K
\end{equation}
Since $f$ is bounded below by $f^*$, we have $f(z_1) - f(z_{T+1}) \leq f(z_1) - f^*$. Dividing by $T \eta \eta_L$ yields:
\begin{equation}
    \frac{1}{T} \sum_{t=1}^T \Phi_t \leq \frac{f(z_1) - f^*}{T \eta \eta_L} + \frac{K}{\eta \eta_L}
\end{equation}

\subsection*{Rate Derivation (With Constants)}
Observe the structure of the constant $K$. Expanding $\Delta_z^2$, we see that $K$ contains terms of order $O(\eta^2)$ and $O(\eta \sqrt{1-\beta_2})$. 
Specifically, $K = c_1 \eta^2 + c_2 \eta \sqrt{1-\beta_2}$, where $c_1$ and $c_2$ depend strictly on $L, \beta_1, \eta_U, C, G$, and $\epsilon$.
Substitute this into our final bound:
\begin{equation}
    \frac{1}{T} \sum_{t=1}^T \Phi_t \leq \frac{f(z_1) - f^*}{T \eta \eta_L} + \frac{c_1 \eta}{\eta_L} + \frac{c_2 \sqrt{1-\beta_2}}{\eta_L}
\end{equation}
To minimize this bound with respect to $T$, we set the learning rate $\eta = \frac{1}{\sqrt{T}}$ and schedule the second moment parameter such that $1-\beta_2 = \frac{1}{T}$.
\begin{align}
    \frac{1}{T} \sum_{t=1}^T \Phi_t &\leq \frac{f(z_1) - f^*}{\sqrt{T} \eta_L} + \frac{c_1}{\eta_L \sqrt{T}} + \frac{c_2}{\eta_L \sqrt{T}} \nonumber \\
    &= O\left(\frac{1}{\sqrt{T}}\right)
\end{align}
This proves that the algorithm converges to a stationary point at a rate of $O(1/\sqrt{T})$.

\subsection*{Special Cases}
\begin{itemize}
    \item \textbf{Clipped RMSProp ($\beta_1 = 0$):} If momentum is disabled, the virtual sequence $z_t$ perfectly aligns with $x_t$ ($\|x_t - z_t\| = 0$). The error terms scaling with $\frac{\beta_1}{1-\beta_1}$ vanish entirely, simplifying $K$ and strictly enforcing monotonic expected descent without virtual sequence artifacts.
    \item \textbf{Absence of Clipping ($C \to \infty$):} As $C \to \infty$, the algorithm reverts to standard Adam. However, the theoretical bound diverges because $\eta_L = \frac{1}{C+\epsilon} \to 0$, rendering the right-hand side of the inequality infinite. This perfectly aligns with known theoretical results (e.g., Reddi et al., 2018) demonstrating that standard Adam can fail to converge for non-convex functions without explicit gradient bounds or clipping.
\end{itemize}

\end{document}